\documentclass[12pt]{article}

\input{../../../_assets/preambulo.tex}

\begin{document}

    \input{../../../_assets/portada}
    \portadaExamenFotoDif{../../dueA4.jpg}{GKI\\Exam II}{GKI. Exam II}{MidnightBlue}{-9}{30}{2026}{Arturo Olivares Martos}

    \begin{description}
        \item[Subject] Fundamentals of Artificial Intelligence.
        \item[Academic Year] Summer Semester 2023.
        % \item[Degree] ---.
        % \item[Group] ---.
        \item[Lecturer] Gérald Kämmerer.
        \item[Description] Final Exam.
        \item[Date] 01.08.2023.
        \item[Duration] 90 minutes.
    \end{description}
    \newpage

    % ------------------------------------
    % PART 1: SHORT QUESTIONS
    % ------------------------------------

    \section*{Part 1: Short Questions (54 pts.)}

    \begin{ejercicio}[Unsupervised Learning (3 pts.)]
        What is unsupervised learning? Mark whether each of the following methods belongs to unsupervised learning or not.

        \begin{table}[h!]
        \centering
        \begin{tabular}{lcc}
        \hline
        \textbf{Method} & \textbf{Right} & \textbf{Wrong} \\ \hline
        PCA & $\square$ & $\square$ \\
        Hill Climbing & $\square$ & $\square$ \\
        kNN & $\square$ & $\square$ \\
        kMeans & $\square$ & $\square$ \\
        t-SNE & $\square$ & $\square$ \\
        GANs & $\square$ & $\square$ \\ \hline
        \end{tabular}
        \caption{Classification of Unsupervised Learning Methods}
        \label{tab:unsupervised_methods}
        \end{table}
    \end{ejercicio}

    \begin{ejercicio}[KDD Process (2 pts.)]
        What are the steps involved in the Knowledge Discovery in Databases (KDD) process?
    \end{ejercicio}

    \begin{ejercicio}[Date Conversion Issues (1 pt.)]
        Indicate a possible problem when converting the date $08/01/23$.
    \end{ejercicio}

    \begin{ejercicio}[Basic Statistics (3 pts.)]
        Given the numbers $A = \{3, 4, 5, 10, 11, 1, 8\}$, compute the median, mean, and variance.
    \end{ejercicio}

    \begin{ejercicio}[Data Aggregation Problems (2 pts.)]
        What problems can occur when performing data aggregation? Mark whether the statement is right or wrong.

        \begin{table}[h!]
        \centering
        \begin{tabular}{lcc}
        \hline
        \textbf{Problem} & \textbf{Right} & \textbf{Wrong} \\ \hline
        Information gain & $\square$ & $\square$ \\
        Information loss & $\square$ & $\square$ \\
        Data reduction & $\square$ & $\square$ \\
        Data separation & $\square$ & $\square$ \\ \hline
        \end{tabular}
        \caption{Data Aggregation Issues}
        \label{tab:aggregation_issues}
        \end{table}
    \end{ejercicio}

    \begin{ejercicio}[Curse of Dimensionality (2 pts.)]
        Where does the effect of the ``Curse of Dimensionality'' occur?

        \begin{table}[h!]
        \centering
        \begin{tabular}{lcc}
        \hline
        \textbf{Method} & \textbf{Right} & \textbf{Wrong} \\ \hline
        High-dimensional search & $\square$ & $\square$ \\
        Data aggregation & $\square$ & $\square$ \\
        Optimize a neural network & $\square$ & $\square$ \\
        Execute a binary classifier & $\square$ & $\square$ \\ \hline
        \end{tabular}
        \caption{Curse of Dimensionality Contexts}
        \label{tab:curse_dimensionality}
        \end{table}
    \end{ejercicio}

    \begin{ejercicio}[Search Methods (1 pt.)]
        For which search method is Depth-Limited Search a modification?
    \end{ejercicio}

    \begin{ejercicio}[Topological Search (5 pts.)]
        Give an example to explain Topological Search.
    \end{ejercicio}

    \begin{ejercicio}[Genetic Algorithms - Mutation \& Crossover (3 pts.)]
        Given two genes of length $10$ with bases $\{0, 1, 2\}$:
        \begin{itemize}
            \item Gene 1: $0101211221$
            \item Gene 2: $2111111001$
        \end{itemize}
        Apply mutation and a crossover type to these genes starting with Gene 1.
    \end{ejercicio}

    \begin{ejercicio}[Uniform Cost Search (4 pts.)]
        Explain Uniform Cost Search using the graph in Figure~\ref{fig:uniform_cost_graph}. Go through the different steps.

        \begin{figure}[h!]
        \centering
        \begin{tikzpicture}[
            node/.style={circle, draw, fill=black, text=white, minimum size=1.2cm, align=center, font=\small\sffamily\bfseries},
            edge/.style={-Stealth, thick}
        ]
        \node[node] (A) at (0,0) {Start A};
        \node[node] (B) at (3,2) {B};
        \node[node] (C) at (3,-2) {C};
        \node[node] (D) at (6,3.5) {D};
        \node[node] (E) at (6,0) {E};
        \node[node] (F) at (6,-3.5) {F};
        \node[node] (G) at (9,-2) {G};
        \node[node] (H) at (9,2) {Goal H};

        \draw[edge] (A) -- node[above left, text=black] {2} (B);
        \draw[edge] (A) -- node[below left, text=black] {1} (C);
        \draw[edge] (B) -- node[above left, text=black] {7} (D);
        \draw[edge] (B) -- node[above right, text=black] {3} (E);
        \draw[edge] (C) -- node[below right, text=black] {5} (E);
        \draw[edge] (C) -- node[below left, text=black] {2} (F);
        \draw[edge] (D) -- node[above right, text=black] {1} (H);
        \draw[edge] (E) -- node[above left, text=black] {3} (H);
        \draw[edge, -Stealth, Stealth-] (E) -- node[above, text=black] {4} (G);
        \draw[edge] (F) -- node[below right, text=black] {8} (G);
        \end{tikzpicture}
        \caption{Uniform Cost Graph}
        \label{fig:uniform_cost_graph}
        \end{figure}
    \end{ejercicio}

    \begin{ejercicio}[Roulette Selection (2 pts.)]
        What is Roulette Selection in genetic algorithms?
    \end{ejercicio}

    \begin{ejercicio}[Underfitting vs. Overfitting (2 pts.)]
        Explain what is meant by underfitting and overfitting in machine learning.
    \end{ejercicio}

    \begin{ejercicio}[Confusion Matrix Elements (6 pts.)]
        Fill in the missing entries correctly in Table~\ref{tab:confusion_matrix} and explain the elements of the confusion matrix.

        \begin{table}[h!]
        \centering
        \begin{tabular}{|l|c|c|c|}
        \hline
        \textbf{Predicted / Actual} & \textbf{Actual Yes} & \textbf{Actual No} & \textbf{Actual Total} \\ \hline
        \textbf{Predicted Yes} & 45\% & 15\% & \\ \hline
        \textbf{Predicted No} & & 15\% & 40\% \\ \hline
        \textbf{Predicted Total} & 70\% & & \\ \hline
        \end{tabular}
        \caption{Confusion Matrix}
        \label{tab:confusion_matrix}
        \end{table}
    \end{ejercicio}

    \begin{ejercicio}[F1-Score Definition (1 pt.)]
        How is the F1-Score defined?
    \end{ejercicio}

    \begin{ejercicio}[Resampling Methods (2 pts.)]
        Name one resampling method and explain how it works.
    \end{ejercicio}

    \begin{ejercicio}[Linear vs. Ridge vs. Lasso (3 pts.)]
        List the differences between Linear Regression, Ridge Regression, and Lasso Regression.
    \end{ejercicio}

    \begin{ejercicio}[Principal Component Analysis (1 pt.)]
        What happens during Principal Component Analysis (PCA)?
    \end{ejercicio}

    \begin{ejercicio}[Top-Down Approach (2 pts.)]
        What is the top-down approach in hierarchical clustering or decision trees?
    \end{ejercicio}

    \begin{ejercicio}[Linkage Types in Hierarchical Clustering (4 pts.)]
        Name two types of linkage in hierarchical clustering and state the mutual advantages and disadvantages of each.
    \end{ejercicio}

    \begin{ejercicio}[Activation Function Definition (1 pt.)]
        Name an activation function and give its mathematical definition.
    \end{ejercicio}

    \begin{ejercicio}[Feedback Connections in Neural Networks (3 pts.)]
        Draw and label a lateral feedback connection, a direct feedback connection, and an indirect feedback connection in the network given in Figure~\ref{fig:forward_network}.

        \begin{figure}[h!]
        \centering
        \begin{tikzpicture}[
            node/.style={circle, draw, fill=black, text=white, minimum size=1.2cm, align=center, font=\small\sffamily\bfseries},
            edge/.style={-Stealth, thick}
        ]
        \node[node] (A) at (0,0) {A};
        \node[node] (B) at (3,1.5) {B};
        \node[node] (C) at (3,-1.5) {C};
        \node[node] (D) at (6,3) {D};
        \node[node] (E) at (6,0) {E};
        \node[node] (F) at (6,-3) {F};
        \node[node] (G) at (9,1.5) {G};
        \node[node] (H) at (9,-1.5) {H};
        \node[node] (I) at (12,0) {I};

        \draw[edge] (A) -- (B);
        \draw[edge] (A) -- (C);
        \draw[edge] (B) -- (D);
        \draw[edge] (B) -- (E);
        \draw[edge] (C) -- (E);
        \draw[edge] (C) -- (F);
        \draw[edge] (D) -- (G);
        \draw[edge] (E) -- (G);
        \draw[edge] (E) -- (H);
        \draw[edge] (F) -- (H);
        \draw[edge] (G) -- (I);
        \draw[edge] (H) -- (I);
        \end{tikzpicture}
        \caption{Forward-Feeding Network}
        \label{fig:forward_network}
        \end{figure}
    \end{ejercicio}

    \begin{ejercicio}[Stemming (1 pt.)]
        What is stemming in Natural Language Processing (NLP)?
    \end{ejercicio}

    \newpage

    % ------------------------------------
    % PART 2: KNN
    % ------------------------------------

    \section*{Part 2: kNN Classification (15 pts.)}

    The precipitation dataset (Table~\ref{tab:weather_data}) is provided. The goal is to classify the two days from Table~\ref{tab:data_to_classify} as rainy or non-rainy days.

    \begin{table}[h!]
    \centering
    \begin{tabular}{|c|c|c|c|c|}
    \hline
    \textbf{Day} & \textbf{Max. Temperature} & \textbf{Cloudiness} & \textbf{Rain Level} & \textbf{Rainy Day} \\ \hline
    1 & $20.1^{\circ}\text{C}$ & 30\% & 0.30 mm & Yes \\ \hline
    2 & $29.8^{\circ}\text{C}$ & 80\% & 0.28 mm & Yes \\ \hline
    3 & $21.2^{\circ}\text{C}$ & 25\% & 0.00 mm & No \\ \hline
    4 & $25.0^{\circ}\text{C}$ & 20\% & 0.01 mm & No \\ \hline
    5 & $22.5^{\circ}\text{C}$ & 85\% & 0.35 mm & Yes \\ \hline
    6 & $19.8^{\circ}\text{C}$ & 30\% & 0.05 mm & No \\ \hline
    7 & $20.2^{\circ}\text{C}$ & 80\% & 0.31 mm & Yes \\ \hline
    \end{tabular}
    \caption{Weather Data for kNN Classification}
    \label{tab:weather_data}
    \end{table}

    \begin{ejercicio}[Feature Types (2 pts.)]
        Which features in Table~\ref{tab:weather_data} are categorical and which are nominal?
    \end{ejercicio}

    \begin{ejercicio}[Distance Metrics (5 pts.)]
        Give two different distance metrics and provide their mathematical formulas. Explain the formulas and the differences between them.
    \end{ejercicio}

    \begin{ejercicio}[kNN Execution (8 pts.)]
        Classify the two days given in Table~\ref{tab:data_to_classify} using the kNN algorithm with $k=4$.

        \begin{table}[h!]
        \centering
        \begin{tabular}{|c|c|c|c|}
        \hline
        \textbf{Day} & \textbf{Max. Temperature} & \textbf{Cloudiness} & \textbf{Rain Level} \\ \hline
        1 & $21.3^{\circ}\text{C}$ & 45\% & 0.25 mm \\ \hline
        2 & $18.3^{\circ}\text{C}$ & 25\% & 0.00 mm \\ \hline
        \end{tabular}
        \caption{Data to Classify}
        \label{tab:data_to_classify}
        \end{table}

        \textbf{Hint:} Use the following modified distance metric:
        \[ d(\text{Day } 1, \text{Day } 2) = |\text{Difference in Cloudiness } (\%)| + 100 \cdot |\text{Difference in Rain Level } (\text{mm})| \]

        \textbf{Example:}
        \[ d(\text{Day } 1, \text{Day } 2) = |30 - 80| + 100 \cdot |0.30 - 0.28| = 50 + 100 \cdot 0.02 = 50 + 2 = 52 \]

        Specify the distance ranking explicitly for both days and state the final classification decision.
    \end{ejercicio}

    \newpage

    % ------------------------------------
    % PART 3: DECISION TREE & RANDOM FOREST
    % ------------------------------------

    \section*{Part 3: Decision Tree \& Random Forest (20 pts.)}

    \begin{ejercicio}[ID3 Algorithm Steps (4 pts.)]
        Name and explain the steps of the ID3 algorithm.
    \end{ejercicio}

    \begin{ejercicio}[ID3 vs. C4.5 (2 pts.)]
        Name one difference between the ID3 and C4.5 algorithms.
    \end{ejercicio}

    \begin{ejercicio}[Gini Impurity Calculation (3 pts.)]
        Consider the dataset with 20 points given in Table~\ref{tab:fruit_dataset} and illustrated in Figure~\ref{fig:fruit_scatter_plot}. Features $F_1$ and $F_2$ correspond to $Feature1$ and $Feature2$. Calculate the individual Gini impurities and the total Gini impurity for the split at $F_2 = 4$ based on the target class \textit{Type} (Pear/Apple).

        \begin{figure}[h!]
        \centering
        \begin{tikzpicture}[scale=0.9, font=\sffamily\scriptsize]
            % Grid
            \draw[step=0.5, gray!20, thin] (0,0) grid (9,6.5);
            \draw[step=1.0, gray!40, stroke] (0,0) grid (9,6.5);
            
            % Axes
            \draw[-Stealth, thick] (0,0) -- (9.2,0) node[below] {Feature1};
            \draw[-Stealth, thick] (0,0) -- (0,6.7) node[above] {Feature2};
            
            % Ticks
            \foreach \x in {1,...,9} \draw (\x,0.05) -- (\x,-0.05) node[below] {\x};
            \foreach \y in {1,...,6} \draw (0.05,\y) -- (-0.05,\y) node[left] {\y};
            
            % Data Points (Pears: Red Squares, Apples: Blue Circles)
            % Pears
            \node[circle, fill=red!80!black, inner sep=1.5pt, label={above:\tiny 1}] at (2.00, 5.75) {};
            \node[circle, fill=red!80!black, inner sep=1.5pt, label={above:\tiny 2}] at (1.25, 3.75) {};
            \node[circle, fill=red!80!black, inner sep=1.5pt, label={above:\tiny 3}] at (3.50, 3.75) {};
            \node[circle, fill=red!80!black, inner sep=1.5pt, label={above:\tiny 4}] at (4.00, 5.75) {};
            \node[circle, fill=red!80!black, inner sep=1.5pt, label={above:\tiny 5}] at (1.00, 2.00) {};
            \node[circle, fill=red!80!black, inner sep=1.5pt, label={above:\tiny 6}] at (2.75, 2.25) {};
            \node[circle, fill=red!80!black, inner sep=1.5pt, label={above:\tiny 7}] at (1.25, 1.25) {};
            \node[circle, fill=red!80!black, inner sep=1.5pt, label={above:\tiny 11}] at (5.75, 5.50) {};
            \node[circle, fill=red!80!black, inner sep=1.5pt, label={above:\tiny 13}] at (7.75, 5.50) {};
            \node[circle, fill=red!80!black, inner sep=1.5pt, label={above:\tiny 15}] at (8.00, 3.50) {};

            % Apples
            \node[circle, fill=blue!80!black, inner sep=1.5pt, label={above:\tiny 8}] at (3.50, 1.75) {};
            \node[circle, fill=blue!80!black, inner sep=1.5pt, label={above:\tiny 9}] at (4.75, 0.50) {};
            \node[circle, fill=blue!80!black, inner sep=1.5pt, label={above:\tiny 10}] at (5.25, 4.25) {};
            \node[circle, fill=blue!80!black, inner sep=1.5pt, label={above:\tiny 12}] at (5.50, 2.50) {};
            \node[circle, fill=blue!80!black, inner sep=1.5pt, label={above:\tiny 14}] at (8.25, 4.50) {};
            \node[circle, fill=blue!80!black, inner sep=1.5pt, label={above:\tiny 16}] at (7.75, 3.75) {};
            \node[circle, fill=blue!80!black, inner sep=1.5pt, label={above:\tiny 17}] at (6.50, 3.50) {};
            \node[circle, fill=blue!80!black, inner sep=1.5pt, label={above:\tiny 18}] at (6.50, 2.25) {};
            \node[circle, fill=blue!80!black, inner sep=1.5pt, label={above:\tiny 19}] at (6.75, 3.25) {};
            \node[circle, fill=blue!80!black, inner sep=1.5pt, label={above:\tiny 20}] at (7.00, 2.00) {};
            
            % Legend
            \begin{scope}[shift={(6.5, 0.8)}]
                \draw[fill=white, draw=gray] (-0.2,-0.2) rectangle (2.2, 0.8);
                \node[circle, fill=red!80!black, inner sep=1.5pt, label={right:\tiny Pear}] at (0,0.5) {};
                \node[circle, fill=blue!80!black, inner sep=1.5pt, label={right:\tiny Apple}] at (0,0.1) {};
            \end{scope}
        \end{tikzpicture}
        \caption{Scatter Plot of Feature1 vs Feature2}
        \label{fig:fruit_scatter_plot}
        \end{figure}

        \begin{table}[h!]
        \centering
        \small
        \begin{tabular}{|c|c|c|c|}
        \hline
        \textbf{Index} & $F_1$ & $F_2$ & \textbf{Type} \\ \hline
        1 & 2.00 & 5.75 & Pear \\ \hline
        2 & 1.25 & 3.75 & Pear \\ \hline
        3 & 3.50 & 3.75 & Pear \\ \hline
        4 & 4.00 & 5.75 & Pear \\ \hline
        5 & 1.00 & 2.00 & Pear \\ \hline
        6 & 2.75 & 2.25 & Pear \\ \hline
        7 & 1.25 & 1.25 & Pear \\ \hline
        8 & 3.50 & 1.75 & Apple \\ \hline
        9 & 4.75 & 0.50 & Apple \\ \hline
        10 & 5.25 & 4.25 & Apple \\ \hline
        11 & 5.75 & 5.50 & Pear \\ \hline
        12 & 5.50 & 2.50 & Apple \\ \hline
        13 & 7.75 & 5.50 & Pear \\ \hline
        14 & 8.25 & 4.50 & Apple \\ \hline
        15 & 8.00 & 3.50 & Pear \\ \hline
        16 & 7.75 & 3.75 & Apple \\ \hline
        17 & 6.50 & 3.50 & Apple \\ \hline
        18 & 6.50 & 2.25 & Apple \\ \hline
        19 & 6.75 & 3.25 & Apple \\ \hline
        20 & 7.00 & 2.00 & Apple \\ \hline
        \end{tabular}
        \caption{Fruit Dataset}
        \label{tab:fruit_dataset}
        \end{table}
    \end{ejercicio}

    \begin{ejercicio}[Decision Tree Construction (5 pts.)]
        Table~\ref{tab:feature_splits} provides several candidate splits and their corresponding Gini Impurities. Create a Decision Tree that correctly classifies all points based on this table. Draw it explicitly and specify its depth.

        \begin{table}[h!]
        \centering
        \begin{tabular}{|c|c|c|}
        \hline
        \textbf{Feature} & \textbf{Split} & \textbf{Gini Impurity} \\ \hline
        $F_1$ & $<1$ & 0.487 \\ \hline
        $F_1$ & $\le 2$ & 0.469 \\ \hline
        $F_1$ & $<3$ & 0.313 \\ \hline
        $F_1$ & $<4$ & 0.444 \\ \hline
        $F_1$ & $<5$ & 0.374 \\ \hline
        $F_1$ & $<6$ & 0.487 \\ \hline
        $F_1$ & $\le 7$ & 0.500 \\ \hline
        $F_1$ & $<8$ & 0.421 \\ \hline
        $F_2$ & $<1$ & 0.498 \\ \hline
        $F_2$ & $<2$ & 0.493 \\ \hline
        $F_2$ & $<3$ & 0.479 \\ \hline
        $F_2$ & $<4$ & 0.493 \\ \hline
        $F_2$ & $<5$ & 0.375 \\ \hline
        \end{tabular}
        \caption{Feature Splits with Gini Impurity}
        \label{tab:feature_splits}
        \end{table}
    \end{ejercicio}

    \begin{ejercicio}[Bootstrap Samples (1 pt.)]
        Why are bootstrap samples used when building a Random Forest?
    \end{ejercicio}

    \begin{ejercicio}[Random Forest vs. Decision Tree Classification (5 pts.)]
        Classify the point $(F_1 = 4, F_2 = 5.25)$ using the Random Forest in Figure ~\ref{fig:random_forest_decision_tree}. State the prediction of each individual tree as well as the aggregated output. Then, classify the same point using the large Decision Tree on the right (draw the decision path).

        \begin{figure}[h!]
        \centering
        \resizebox{\textwidth}{!}{%
        \begin{tikzpicture}[
            node/.style={circle, draw=black, fill=black, text=white, minimum size=0.95cm, inner sep=0pt, font=\scriptsize\bfseries, align=center},
            leaf/.style={circle, draw=black, fill=black, text=white, minimum size=0.85cm, inner sep=0pt, font=\scriptsize\bfseries},
            edge from parent/.style={draw, -Stealth, ultra thick},
            tree label/.style={font=\small\bfseries, yshift=0.35cm}
        ]

        % Small Tree 1
        \begin{scope}[shift={(0,0)}]
            \node[node] at (0,0) {$F_2 < 4$}
                [sibling distance=1.6cm, level distance=1.3cm]
                child { node[leaf] {C2} }
                child { node[leaf] {C3} };
        \end{scope}

        % Small Tree 2
        \begin{scope}[shift={(3.5,0)}]
            \node[node] at (0,0) {$F_1 \ge 4$}
                [sibling distance=1.6cm, level distance=1.3cm]
                child { node[leaf] {C2} }
                child { node[leaf] {C3} };
        \end{scope}

        % Small Tree 3
        \begin{scope}[shift={(0,-3.5)}]
            \node[node] at (0,0) {$F_1 < 3$}
                [sibling distance=1.6cm, level distance=1.3cm]
                child { node[leaf] {C1} }
                child { node[leaf] {C2} };
        \end{scope}

        % Small Tree 4
        \begin{scope}[shift={(3.5,-3.5)}]
            \node[node] at (0,0) {$F_2 > 6$}
                [sibling distance=1.6cm, level distance=1.3cm]
                child { node[leaf] {C1} }
                child { node[leaf] {C2} };
        \end{scope}

        % Large Decision Tree
        \begin{scope}[shift={(10,0)}]
            \node[node] at (0,0) {$F_1 > 3$}
                [sibling distance=3.6cm, level distance=1.3cm]
                child { node[node] {$F_2 > 4$}
                    [sibling distance=1.8cm, level distance=1.3cm]
                    child { node[node] {$F_2 < 5$}
                        [sibling distance=1.2cm, level distance=1.3cm]
                        child { node[leaf] {C1} }
                        child { node[leaf] {C2} }
                    }
                    child { node[leaf] {C4} }
                }
                child { node[node] {$F_2 < 3$}
                    [sibling distance=1.8cm, level distance=1.3cm]
                    child { node[leaf] {C4} }
                    child { node[node] {$F_1 < 4$}
                        [sibling distance=1.2cm, level distance=1.3cm]
                        child { node[leaf] {C3} }
                        child { node[leaf] {C5} }
                    }
                };
        \end{scope}

        \end{tikzpicture}
        }
        \caption{Random Forest \& Decision Tree}
        \label{fig:random_forest_decision_tree}
        \end{figure}
    \end{ejercicio}

    \newpage

    % ------------------------------------
    % PART 4: ENCODING & FORWARD PROPAGATION
    % ------------------------------------

    \section*{Part 4: Encoding \& Forward Propagation (5 pts.)}

    \begin{ejercicio}[Ordinal Encoding (2 pts.)]
        Given a table column with the elements \{\text{``OH''}, \text{``OR''}, \text{``FI''}, \text{``NY''}, \text{``WY''}, \text{``WI''}\}, perform ordinal encoding on these expressions and present the output in a table. What happens if you now want to add \text{``CA''} to the list?
    \end{ejercicio}

    \begin{ejercicio}[Forward Propagation Network (2 pts.)]
        Given the small neural network below with input $x = 12$, calculate the output values after applying the activation functions and write them into the corresponding white boxes.

        \begin{figure}[h!]
        \centering
        \resizebox{0.8\textwidth}{!}{%
        \begin{tikzpicture}[
            neuron/.style={circle, draw, fill=black, text=white, minimum size=1.3cm, font=\small\sffamily\bfseries, align=center},
            box/.style={rectangle, draw, minimum size=0.65cm, thick, fill=white},
            edge/.style={-Stealth, thick},
            weight/.style={font=\small, fill=white, inner sep=1.5pt, rounded corners=1pt, sloped}
        ]

        % Nodes Placement
        \node[font=\large\bfseries] (in) at (-1.5,0) {$x = 12$};
        \node[neuron] (Identity) at (1.5,0) {Identity};
        \node[box] (Identity_box) at (2.9,0) {};

        \node[neuron] (R1) at (6, 2.2) {ReLU};
        \node[box] (R1_box) at (6, 3.3) {};

        \node[neuron] (R2) at (6, -2.2) {ReLU};
        \node[box] (R2_box) at (6, -3.3) {};

        \node[neuron] (Sign) at (10.5,0) {Sign};
        \node[box] (Out_box) at (12.2,0) {};

        % Connections & Weights
        \draw[edge] (in) -- (Identity);
        \draw[edge] (Identity_box) -- node[weight, above] {$w_{1,1} = 4$} (R1);
        \draw[edge] (Identity_box) -- node[weight, below] {$w_{1,2} = -2$} (R2);

        \draw[edge] (R1) -- node[weight, above] {$w_{2,1} = 0.5$} (Sign);
        \draw[edge] (R2) -- node[weight, below] {$w_{2,2} = 0.5$} (Sign);
        \draw[edge] (Sign) -- (Out_box);

        \end{tikzpicture}
        }
        \caption{Small Feedforward Neural Network}
        \end{figure}
    \end{ejercicio}

    \begin{ejercicio}[Error Calculation - MAE \& MSE (1 pt.)]
        Suppose the neural network outputs $2$, and the target value should be $4$. Calculate the Mean Absolute Error (MAE) and Mean Squared Error (MSE).
    \end{ejercicio}


    \setcounter{section}{0}
    \setcounter{ejercicio}{0}

    \newpage\newpage
    % ------------------------------------
    % PART 1: SHORT QUESTIONS
    % ------------------------------------

    \section*{Part 1: Short Questions (54 pts.)}

    \begin{ejercicio}[Unsupervised Learning (3 pts.)]
        What is unsupervised learning? Mark whether each of the following methods belongs to unsupervised learning or not.

        \begin{table}[h!]
        \centering
        \begin{tabular}{lcc}
        \hline
        \textbf{Method} & \textbf{Right} & \textbf{Wrong} \\ \hline
        PCA & $\boxtimes$ & $\square$ \\
        Hill Climbing & $\square$ & $\boxtimes$ \\
        kNN & $\square$ & $\boxtimes$ \\
        kMeans & $\boxtimes$ & $\square$ \\
        t-SNE & $\boxtimes$ & $\square$ \\
        GANs & $\boxtimes$ & $\square$ \\ \hline
        \end{tabular}
        \caption{Classification of Unsupervised Learning Methods}
        \label{tab:unsupervised_methods2}
        \end{table}
    \end{ejercicio}

    \begin{ejercicio}[KDD Process (2 pts.)]
        What are the steps involved in the Knowledge Discovery in Databases (KDD) process?
    \end{ejercicio}

    \begin{ejercicio}[Date Conversion Issues (1 pt.)]
        Indicate a possible problem when converting the date $08/01/23$.

        It can either be understood as August 1, 2023, or January 8, 2023, depending on the date format used (MM/DD/YY vs. DD/MM/YY).
    \end{ejercicio}

    \begin{ejercicio}[Basic Statistics (3 pts.)]
        Given the numbers $A = \{3, 4, 5, 10, 11, 1, 8\}$, compute the median, mean, and variance.
        \begin{itemize}
            \item For the median, sort the numbers: $1, 3, 4, 5, 8, 10, 11$. The median is the middle value: $5$.
            \item For the mean, calculate the sum of the numbers and divide by the count:
            \[ \text{Mean} = \frac{3 + 4 + 5 + 10 + 11 + 1 + 8}{7} = \frac{42}{7} = 6. \]
            \item For the variance, the formula is:
            \begin{align*}
                \text{Variance} &= \frac{1}{n} \sum_{i=1}^{n} (x_i - \text{Mean})^2 \\
                &= \frac{1}{7} \left[3^2+2^2+1^2+4^2+5^2+5^2+2^2\right] \\
                &= \frac{84}{7}= 12
            \end{align*}
        \end{itemize}
    \end{ejercicio}

    \begin{ejercicio}[Data Aggregation Problems (2 pts.)]
        What problems can occur when performing data aggregation? Mark whether the statement is right or wrong.

        \begin{table}[h!]
        \centering
        \begin{tabular}{lcc}
        \hline
        \textbf{Problem} & \textbf{Right} & \textbf{Wrong} \\ \hline
        Information gain & $\square$ & $\boxtimes$ \\
        Information loss & $\boxtimes$ & $\square$ \\
        Data reduction & $\boxtimes$ & $\square$ \\
        Data separation & $\square$ & $\boxtimes$ \\ \hline
        \end{tabular}
        \caption{Data Aggregation Issues}
        \label{tab:aggregation_issues2}
        \end{table}
    \end{ejercicio}

    \begin{ejercicio}[Curse of Dimensionality (2 pts.)]
        Where does the effect of the ``Curse of Dimensionality'' occur?

        \begin{table}[h!]
        \centering
        \begin{tabular}{lcc}
        \hline
        \textbf{Method} & \textbf{Right} & \textbf{Wrong} \\ \hline
        High-dimensional search & $\boxtimes$ & $\square$ \\
        Data aggregation & $\square$ & $\boxtimes$ \\
        Optimize a neural network & $\boxtimes$ & $\square$ \\
        Execute a binary classifier & $\square$ & $\boxtimes$ \\ \hline
        \end{tabular}
        \caption{Curse of Dimensionality Contexts}
        \label{tab:curse_dimensionality2}
        \end{table}
    \end{ejercicio}

    \begin{ejercicio}[Search Methods (1 pt.)]
        For which search method is Depth-Limited Search a modification?

        Depth-Limited Search is a modification of Depth-First Search (DFS), where a depth limit is imposed to prevent the search from going too deep into the search tree, which can help avoid infinite loops in cases of cyclic graphs.
    \end{ejercicio}

    \begin{ejercicio}[Topological Search (5 pts.)]
        Give an example to explain Topological Search.

        Topological Search is used in directed acyclic graphs (DAGs) to order the vertices such that for every directed edge from vertex A to vertex B, A comes before B in the ordering. An example could be deciding the order in which a student has to take courses based on prerequisites. If Course A is a prerequisite for Course B, then in the topological order, Course A will appear before Course B.
    \end{ejercicio}

    \begin{ejercicio}[Genetic Algorithms - Mutation \& Crossover (3 pts.)]
        Given two genes of length $10$ with bases $\{0, 1, 2\}$:
        \begin{itemize}
            \item Gene 1: $0101211221$
            \item Gene 2: $2111111001$
        \end{itemize}
        Apply mutation and a crossover type to these genes starting with Gene 1.
        \begin{table}[ht]
            \centering
            \begin{tabular}{r|cccccccccc}
                Parent 1 & 0 & 1 & 0 & 1 & 2 & 1 & 1 & 2 & 2 & 1 \\
                Parent 2 & 2 & 1 & 1 & 1 & 1 & 1 & 1 & 0 & 0 & 1 \\ \hline
                Offspring & 2 & 1 & 0 & 1 & 2 & 1 & 1 & 0 & 0 & 1 \\
                Swap Mutation & 2 & 1 & 0 & 1 & 2 & 1 & 1 & 0 & \textbf{1} & \textbf{0} \\
            \end{tabular}
            \caption{\centering Example of Order-Based Crossover (OX2), where the fixed positions are 1, 2, 4, and 6.}
            \label{tab:ox2_example}
        \end{table}
    \end{ejercicio}

    \begin{ejercicio}[Uniform Cost Search (4 pts.)]
        Explain Uniform Cost Search using the graph in Figure~\ref{fig:uniform_cost_graph2}. Go through the different steps.

        \begin{figure}[h!]
        \centering
        \begin{tikzpicture}[
            node/.style={circle, draw, fill=black, text=white, minimum size=1.2cm, align=center, font=\small\sffamily\bfseries},
            edge/.style={-Stealth, thick}
        ]
        \node[node] (A) at (0,0) {Start A};
        \node[node] (B) at (3,2) {B};
        \node[node] (C) at (3,-2) {C};
        \node[node] (D) at (6,3.5) {D};
        \node[node] (E) at (6,0) {E};
        \node[node] (F) at (6,-3.5) {F};
        \node[node] (G) at (9,-2) {G};
        \node[node] (H) at (9,2) {Goal H};

        \draw[edge] (A) -- node[above left, text=black] {2} (B);
        \draw[edge] (A) -- node[below left, text=black] {1} (C);
        \draw[edge] (B) -- node[above left, text=black] {7} (D);
        \draw[edge] (B) -- node[above right, text=black] {3} (E);
        \draw[edge] (C) -- node[below right, text=black] {5} (E);
        \draw[edge] (C) -- node[below left, text=black] {2} (F);
        \draw[edge] (D) -- node[above right, text=black] {1} (H);
        \draw[edge] (E) -- node[above left, text=black] {3} (H);
        \draw[edge, -Stealth, Stealth-] (E) -- node[above, text=black] {4} (G);
        \draw[edge] (F) -- node[below right, text=black] {8} (G);
        \end{tikzpicture}
        \caption{Uniform Cost Graph}
        \label{fig:uniform_cost_graph2}
        \end{figure}

        A Uniform Cost Search (UCS) is an algorithm used to find the least-cost path from a starting node to a goal node in a weighted graph. It expands the node with the lowest cumulative cost first, ensuring that the path found is optimal.
        \begin{itemize}
            \item Start at node A with a cumulative cost of 0. The priority queue contains: $[(A, 0)]$.
            \item Expand node A. The neighbors are B (cost 2) and C (cost 1). Add them to the queue: $[(C, 1), (B, 2)]$.
            \item Expand node C (lowest cost 1). The neighbors are E (cost 6) and F (cost 3). Add them to the queue: $[(B, 2), (F, 3), (E, 6)]$.
            \item Expand node B (lowest cost 2). The neighbors are D (cost 9) and E (cost 5). Add them to the queue: $[(F, 3), (E, 5), (D, 9)]$.
            \item Expand node F (lowest cost 3). The neighbor is G (cost 11). Add it to the queue: $[(E, 5), (D, 9), (G, 11)]$.
            \item Expand node E (lowest cost 5). The neighbors are H (cost 8). Add it to the queue: $[(H, 8), (D, 9), (G, 11)]$.
            \item Expand node H (lowest cost 8). Since H is the goal node, the search terminates, and the least-cost path from A to H is found with a total cost of 8.
        \end{itemize}
    \end{ejercicio}

    \begin{ejercicio}[Roulette Selection (2 pts.)]
        What is Roulette Selection in genetic algorithms?

        Roulette Selection is a selection method used in genetic algorithms to choose individuals from a population for reproduction based on their fitness. Each individual is assigned a probability of being selected that is proportional to its fitness score. The higher the fitness, the greater the chance of being selected. This is often visualized as a roulette wheel where each individual's segment size corresponds to its fitness, and a random spin determines which individual is chosen for the next generation.
        
        This has however one big disadvantage: if one individual has a very high fitness compared to the others, it can dominate the selection process, leading to premature convergence and loss of genetic diversity in the population. Rank selection is a better alternative in such cases, as it assigns selection probabilities based on the rank of individuals rather than their absolute fitness values, ensuring a more balanced selection process.
    \end{ejercicio}

    \begin{ejercicio}[Underfitting vs. Overfitting (2 pts.)]
        Explain what is meant by underfitting and overfitting in machine learning.

        Underfitting occurs when a machine learning model is too simple to capture the underlying patterns in the training data, resulting in poor performance on both the training and test datasets. It fails to learn the relationships between features and target variables, leading to high bias and low variance.

        Overfitting, on the other hand, happens when a model is too complex and learns not only the underlying patterns but also the noise in the training data. This results in excellent performance on the training dataset but poor generalization to unseen test data. Overfitting is characterized by low bias and high variance, where the model captures random fluctuations rather than the true signal.
    \end{ejercicio}

    \begin{ejercicio}[Confusion Matrix Elements (6 pts.)]
        Fill in the missing entries correctly in Table~\ref{tab:confusion_matrix2} and explain the elements of the confusion matrix.

        \begin{table}[h!]
        \centering
        \begin{tabular}{|l|c|c|c|}
        \hline
        \textbf{Predicted / Actual} & \textbf{Actual Yes} & \textbf{Actual No} & \textbf{Actual Total} \\ \hline
        \textbf{Predicted Yes} & 45\% & 15\% & 60\%\\ \hline
        \textbf{Predicted No} & 25\%& 15\% & 40\% \\ \hline
        \textbf{Predicted Total} & 70\% & 30\% & 100\% \\ \hline
        \end{tabular}
        \caption{Confusion Matrix}
        \label{tab:confusion_matrix2}
        \end{table}

        \begin{itemize}
            \item Predicted Yes / Actual Yes (True Positives): 45\% - The model correctly predicted positive cases.
            \item Predicted Yes / Actual No (False Positives): 15\% - The model incorrectly predicted positive cases when they were actually negative.
            \item Predicted No / Actual Yes (False Negatives): 25\% - The model incorrectly predicted negative cases when they were actually positive.
            \item Predicted No / Actual No (True Negatives): 15\% - The model correctly predicted negative cases.
        \end{itemize}
    \end{ejercicio}

    \begin{ejercicio}[F1-Score Definition (1 pt.)]
        How is the F1-Score defined?

        It is the harmonic mean of precision and recall, calculated as:
        \[ F1 = 2 \cdot \frac{\text{Precision} \cdot \text{Recall}}{\text{Precision} + \text{Recall}} \]
    \end{ejercicio}

    \begin{ejercicio}[Resampling Methods (2 pts.)]
        Name one resampling method and explain how it works.

        One resampling method is k-fold cross-validation. In k-fold cross-validation, the dataset is divided into k equal-sized subsets or ``folds''. The model is trained k times, each time using k-1 folds for training and the remaining fold for validation. This process is repeated until each fold has been used as the validation set once. The performance metrics are then averaged over the k iterations to provide a more robust estimate of the model's generalization ability.
    \end{ejercicio}

    \begin{ejercicio}[Linear vs. Ridge vs. Lasso (3 pts.)]
        List the differences between Linear Regression, Ridge Regression, and Lasso Regression.

        The only difference lies on the loss function used to train the model.
        \begin{itemize}
            \item \ul{Linear Regression}: Uses the Ordinary Least Squares (OLS) loss function, which minimizes the sum of squared residuals between the predicted and actual values.
            \begin{equation*}
                \text{Loss}_{\text{OLS}} = \sum_{i=1}^{n} (y_i - \hat{y}_i)^2
            \end{equation*}

            \item \ul{Ridge Regression}: Adds an L2 regularization term to the OLS loss function, which penalizes large coefficients to prevent overfitting. The loss function is:
            \begin{equation*}
                \text{Loss}_{\text{Ridge}} = \sum_{i=1}^{n} (y_i - \hat{y}_i)^2 + \lambda \sum_{j=1}^{p} \beta_j^2
            \end{equation*}

            \item \ul{Lasso Regression}: Adds an L1 regularization term to the OLS loss function, which can shrink some coefficients to zero, effectively performing feature selection. The loss function is:
            \begin{equation*}
                \text{Loss}_{\text{Lasso}} = \sum_{i=1}^{n} (y_i - \hat{y}_i)^2 + \lambda \sum_{j=1}^{p} |\beta_j|
            \end{equation*}
        \end{itemize}
    \end{ejercicio}

    \begin{ejercicio}[Principal Component Analysis (1 pt.)]
        What happens during Principal Component Analysis (PCA)?

        During Principal Component Analysis (PCA), the original dataset is transformed into a new set of orthogonal axes called principal components. These components are linear combinations of the original features and are ordered by the amount of variance they capture from the data. The first principal component captures the maximum variance, followed by the second, and so on. PCA reduces the dimensionality of the data while retaining as much variability as possible, which can help in visualization, noise reduction, and improving model performance.
    \end{ejercicio}

    \begin{ejercicio}[Top-Down Approach (2 pts.)]
        What is the top-down approach in hierarchical clustering or decision trees?

        The top-down approach, also known as divisive clustering, starts with all data points in a single cluster and recursively splits them into smaller clusters.

        In decision trees, the top-down approach begins with the entire dataset at the root node and recursively partitions it into subsets based on feature values, creating branches and nodes until a stopping criterion is met (e.g., maximum depth, minimum samples per leaf, or pure nodes). This method allows for a hierarchical representation of the data, where each split aims to maximize the separation of classes or minimize impurity.
    \end{ejercicio}

    \begin{ejercicio}[Linkage Types in Hierarchical Clustering (4 pts.)]
        Name two types of linkage in hierarchical clustering and state the mutual advantages and disadvantages of each.

        \begin{itemize}
            \item \ul{Single Linkage}: This method defines the distance between two clusters as the minimum distance between any single pair of points in the two clusters.
            \begin{equation*}
                d_{\text{single}}(C_1, C_2) = \min_{x \in C_1, y \in C_2} d(x, y)
            \end{equation*}

            It is sensitive to noise and outliers, which can lead to the ``chaining effect,'' where clusters may be merged based on a single close point, resulting in elongated and less compact clusters. However, it can capture non-elliptical shapes and is computationally efficient.

            \item \ul{Complete Linkage}: This method defines the distance between two clusters as the maximum distance between any single pair of points in the two clusters.
            \begin{equation*}
                d_{\text{complete}}(C_1, C_2) = \max_{x \in C_1, y \in C_2} d(x, y)
            \end{equation*}
            
            It tends to produce more compact and spherical clusters, making it less sensitive to noise and outliers compared to single linkage. However, it may fail to capture elongated or irregularly shaped clusters and can be computationally more intensive.
        \end{itemize}
    \end{ejercicio}

    \begin{ejercicio}[Activation Function Definition (1 pt.)]
        Name an activation function and give its mathematical definition.

        The Rectified Linear Unit (ReLU) is a widely used activation function in neural networks. It is defined mathematically as:
        \[
        \text{ReLU}(x) = \max(0, x)
        \]
        This function outputs the input directly if it is positive; otherwise, it outputs zero. ReLU introduces non-linearity into the model, allowing it to learn complex patterns, and helps mitigate the vanishing gradient problem during training.
    \end{ejercicio}

    \begin{ejercicio}[Feedback Connections in Neural Networks (3 pts.)]
        Draw and label a lateral feedback connection, a direct feedback connection, and an indirect feedback connection in the network given in Figure~\ref{fig:forward_network2}.

        \begin{figure}[h!]
        \centering
        \begin{tikzpicture}[
            node/.style={circle, draw, fill=black, text=white, minimum size=1.2cm, align=center, font=\small\sffamily\bfseries},
            edge/.style={-Stealth, thick},feedback/.style={-Stealth, thick, red, dashed}, % Red dashed style for added connections
            lbl/.style={font=\scriptsize\bfseries, red, fill=white, inner sep=1.5pt, rounded corners=1pt}
        ]
        \node[node] (A) at (0,0) {A};
        \node[node] (B) at (3,1.5) {B};
        \node[node] (C) at (3,-1.5) {C};
        \node[node] (D) at (6,3) {D};
        \node[node] (E) at (6,0) {E};
        \node[node] (F) at (6,-3) {F};
        \node[node] (G) at (9,1.5) {G};
        \node[node] (H) at (9,-1.5) {H};
        \node[node] (I) at (12,0) {I};

        \draw[edge] (A) -- (B);
        \draw[edge] (A) -- (C);
        \draw[edge] (B) -- (D);
        \draw[edge] (B) -- (E);
        \draw[edge] (C) -- (E);
        \draw[edge] (C) -- (F);
        \draw[edge] (D) -- (G);
        \draw[edge] (E) -- (G);
        \draw[edge] (E) -- (H);
        \draw[edge] (F) -- (H);
        \draw[edge] (G) -- (I);
        \draw[edge] (H) -- (I);

        % =========================================================
        % --- ADDED CONNECTIONS (DRAWN & LABELED) ---
        % =========================================================

        % 1. Direct Feedback Connection: E -> E (Self-loop)
        \draw[feedback] (E.south east) .. controls +(0.8, -1.2) and +(-1.2, -0.8) .. (E.south)
            node[lbl, pos=0.5, anchor=north] {Direct Feedback};

        % 2. Indirect Feedback Connection: G -> B (Later layer back to earlier layer)
        \draw[feedback] (G.north) .. controls +(-2.5, 2.5) and +(2.5, 2.5) .. (B.north)
            node[lbl, pos=0.5, anchor=south] {Indirect Feedback};

        % 3. Lateral Connection: D -> E (Same layer connection)
        \draw[feedback] (D) -- (E)
            node[lbl, pos=0.5, anchor=east] {Lateral};
        \end{tikzpicture}
        \caption{Forward-Feeding Network}
        \label{fig:forward_network2}
        \end{figure}
    \end{ejercicio}

    \begin{ejercicio}[Stemming (1 pt.)]
        What is stemming in Natural Language Processing (NLP)?

        Stemming is a text normalization technique in Natural Language Processing (NLP) that reduces words to their base or root form, known as the stem. The process involves removing suffixes and prefixes from words to group together different forms of the same word, which helps in reducing the dimensionality of the text data and improving the performance of NLP models. For example, the words ``running,'' or ``runner,'' can be reduced to the stem ``run.'' Opposite to lemmatization, stemming does not guarantee that the resulting stem is a valid word in the language; it may produce non-dictionary forms.
    \end{ejercicio}

    \newpage

    % ------------------------------------
    % PART 2: KNN
    % ------------------------------------

    \section*{Part 2: kNN Classification (15 pts.)}

    The precipitation dataset (Table~\ref{tab:weather_data2}) is provided. The goal is to classify the two days from Table~\ref{tab:data_to_classify2} as rainy or non-rainy days.

    \begin{table}[h!]
    \centering
    \begin{tabular}{|c|c|c|c|c|}
    \hline
    \textbf{Day} & \textbf{Max. Temperature} & \textbf{Cloudiness} & \textbf{Rain Level} & \textbf{Rainy Day} \\ \hline
    1 & $20.1^{\circ}\text{C}$ & 30\% & 0.30 mm & Yes \\ \hline
    2 & $29.8^{\circ}\text{C}$ & 80\% & 0.28 mm & Yes \\ \hline
    3 & $21.2^{\circ}\text{C}$ & 25\% & 0.00 mm & No \\ \hline
    4 & $25.0^{\circ}\text{C}$ & 20\% & 0.01 mm & No \\ \hline
    5 & $22.5^{\circ}\text{C}$ & 85\% & 0.35 mm & Yes \\ \hline
    6 & $19.8^{\circ}\text{C}$ & 30\% & 0.05 mm & No \\ \hline
    7 & $20.2^{\circ}\text{C}$ & 80\% & 0.31 mm & Yes \\ \hline
    \end{tabular}
    \caption{Weather Data for kNN Classification}
    \label{tab:weather_data2}
    \end{table}

    \begin{ejercicio}[Feature Types (2 pts.)]
        Which features in Table~\ref{tab:weather_data2} are categorical and which are nominal?

        Among all the features in Table~\ref{tab:weather_data2}, the \textbf{Rainy Day} feature is categorical and nominal, as it represents a binary classification (Yes/No) without any inherent order. The other features (\textbf{Max. Temperature}, \textbf{Cloudiness}, and \textbf{Rain Level}) are numerical and continuous, representing measurable quantities.
    \end{ejercicio}

    \begin{ejercicio}[Distance Metrics (5 pts.)]
        Give two different distance metrics and provide their mathematical formulas. Explain the formulas and the differences between them.

        Two commonly used distance metrics are the Euclidean distance and the Manhattan distance.
        \begin{itemize}
            \item \textbf{Euclidean Distance:} This metric calculates the straight-line distance between two points in a multi-dimensional space. The formula for Euclidean distance between two points $P = (p_1, p_2, ..., p_n)$ and $Q = (q_1, q_2, ..., q_n)$ is:
            \[
            d_{\text{Euclidean}}(P, Q) = \sqrt{\sum_{i=1}^{n} (p_i - q_i)^2}
            \]
            It is sensitive to the scale of the features and can be influenced by outliers.

            \item \textbf{Manhattan Distance:} Also known as the L1 distance or taxicab distance, this metric calculates the sum of the absolute differences between the coordinates of two points. The formula for Manhattan distance is:
            \[
            d_{\text{Manhattan}}(P, Q) = \sum_{i=1}^{n} |p_i - q_i|
            \]
            It is less sensitive to outliers compared to Euclidean distance and is often used when the data has a grid-like structure.
        \end{itemize}
        The main difference between the two metrics is that Euclidean distance measures the shortest path (straight line) between points, while Manhattan distance measures the total distance traveled along axes at right angles (like navigating a grid). This can lead to different distance calculations and neighbor selections in kNN, especially in high-dimensional spaces or when features have different scales.
    \end{ejercicio}

    \begin{ejercicio}[kNN Execution (8 pts.)]
        Classify the two days given in Table~\ref{tab:data_to_classify2} using the kNN algorithm with $k=4$.

        \begin{table}[h!]
        \centering
        \begin{tabular}{|c|c|c|c|}
        \hline
        \textbf{Day} & \textbf{Max. Temperature} & \textbf{Cloudiness} & \textbf{Rain Level} \\ \hline
        1 & $21.3^{\circ}\text{C}$ & 45\% & 0.25 mm \\ \hline
        2 & $18.3^{\circ}\text{C}$ & 25\% & 0.00 mm \\ \hline
        \end{tabular}
        \caption{Data to Classify}
        \label{tab:data_to_classify2}
        \end{table}

        \textbf{Hint:} Use the following modified distance metric:
        \[ d(\text{Day } 1, \text{Day } 2) = |\text{Difference in Cloudiness } (\%)| + 100 \cdot |\text{Difference in Rain Level } (\text{mm})| \]

        \textbf{Example:}
        \[ d(\text{Day } 1, \text{Day } 2) = |30 - 80| + 100 \cdot |0.30 - 0.28| = 50 + 100 \cdot 0.02 = 50 + 2 = 52 \]

        Specify the distance ranking explicitly for both days and state the final classification decision.\\
        
        The distance calculations for Day 1 and Day 2 in Table~\ref{tab:data_to_classify2} against the dataset in Table~\ref{tab:weather_data2} are as follows:
        \begin{itemize}
            \item For Day 1 ($21.3^{\circ}\text{C}$, $45\%$, $0.25$ mm):
            \begin{itemize}
                \item Distance to Day 1: $|30 - 45| + 100 \cdot |0.30 - 0.25| = 15 + 5 = 20$
                \item Distance to Day 2: $|80 - 45| + 100 \cdot |0.28 - 0.25| = 35 + 3 = 38$
                \item Distance to Day 3: $|25 - 45| + 100 \cdot |0.00 - 0.25| = 20 + 25 = 45$
                \item Distance to Day 4: $|20 - 45| + 100 \cdot |0.01 - 0.25| = 25 + 24 = 49$
                \item Distance to Day 5: $|85 - 45| + 100 \cdot |0.35 - 0.25| = 40 + 10 = 50$
                \item Distance to Day 6: $|30 - 45| + 100 \cdot |0.05 - 0.25| = 15 + 20 = 35$
                \item Distance to Day 7: $|80 - 45| + 100 \cdot |0.31 - 0.25| = 35 + 6 = 41$
            \end{itemize}

            Therefore, the sorted distances for Day 1 are shown in Table~\ref{tab:day1_distances}.
            \begin{table}[h!]
            \centering
            \begin{tabular}{|c|c|c|}
            \hline
            \textbf{Day} & \textbf{Distance} & \textbf{Rainy Day} \\ \hline
            1 & 20 & Yes \\ 
            6 & 35 & No \\ 
            2 & 38 & Yes \\ 
            7 & 41 & Yes \\ \hline
            3 & 45 & No \\ 
            4 & 49 & No \\ 
            5 & 50 & Yes \\ \hline
            \end{tabular}
            \caption{Distance Ranking for Day 1}
            \label{tab:day1_distances}
            \end{table}
            Since only the first $k=4$ rows are considered, the majority of the nearest neighbors are rainy days (3 out of 4), so Day 1 is classified as a rainy day.

            \item For Day 2 ($18.3^{\circ}\text{C}$, $25\%$, $0.00$ mm):
            \begin{itemize}
                \item Distance to Day1: $|30 - 25| +100\cdot |0.30-0.00|=5+30=35$
                \item Distance to Day2: $|80 - 25| +100\cdot |0.28-0.00|=55+28=83$
                \item Distance to Day3: $|25 - 25| +100\cdot |0.00-0.00|=0+0=0$
                \item Distance to Day4: $|20 - 25| +100\cdot |0.01-0.00|=5+1=6$
                \item Distance to Day5: $|85 - 25| +100\cdot |0.35-0.00|=60+35=95$
                \item Distance to Day6: $|30 - 25| +100\cdot |0.05-0.00|=5+5=10$
                \item Distance to Day7: $|80 - 25| +100\cdot |0.31-0.00|=55+31=86$
            \end{itemize}
            Therefore, the sorted distances for Day 2 are shown in Table~\ref{tab:day2_distances}.
            \begin{table}[h!]
            \centering
            \begin{tabular}{|c|c|c|}
            \hline
            \textbf{Day} & \textbf{Distance} & \textbf{Rainy Day} \\ \hline
            3 & 0 & No \\ 
            4 & 6 & No \\ 
            6 & 10 & No \\ 
            1 & 35 & Yes \\ \hline
            2 & 83 & Yes \\ 
            7 & 86 & Yes \\ 
            5 & 95 & Yes \\ \hline
            \end{tabular}
            \caption{Distance Ranking for Day 2}
            \label{tab:day2_distances}
            \end{table}
            Since only the first $k=4$ rows are considered, the majority of the nearest neighbors are non-rainy days (3 out of 4), so Day 2 is classified as a non-rainy day.
        \end{itemize}
    \end{ejercicio}

    \newpage

    % ------------------------------------
    % PART 3: DECISION TREE & RANDOM FOREST
    % ------------------------------------

    \section*{Part 3: Decision Tree \& Random Forest (20 pts.)}

    \begin{ejercicio}[ID3 Algorithm Steps (4 pts.)]
        Name and explain the steps of the ID3 algorithm.

        The ID3 (Iterative Dichotomiser 3) algorithm is a decision tree learning algorithm that constructs a decision tree by employing a top-down, greedy approach. The steps of the ID3 algorithm are as follows:
        \begin{enumerate}
            \item \textbf{Select the Best Attribute:} For each node, for each attribute, calculate the entropy of the dataset resulting from splitting the data based on that attribute. The attribute with the lowest entropy is selected as the best attribute to split the data at that node.
            \item \textbf{Create a Decision Node:} Create a decision node in the tree that corresponds to the selected attribute. This node will have branches for each possible value of the attribute.
            \item \textbf{Split the Dataset:} Partition the dataset into subsets based on the values of the selected attribute. Each subset corresponds to a branch of the decision node.
            \item \textbf{Repeat Recursively:} For each subset, repeat the process of selecting the best attribute and creating decision nodes until one of the following stopping criteria is met:
            \begin{itemize}
                \item All instances in the subset belong to the same class (pure node).
                \item There are no remaining attributes to split on.
                \item A predefined stopping condition is reached (e.g., maximum depth, minimum samples per leaf).
            \end{itemize}
            \item \textbf{Assign Class Labels:} Once a stopping criterion is met, assign the majority class label of the instances in the subset to the leaf node.
        \end{enumerate}
    \end{ejercicio}

    \begin{ejercicio}[ID3 vs. C4.5 (2 pts.)]
        Name one difference between the ID3 and C4.5 algorithms.


        One main difference between the ID3 and C4.5 algorithms is that C4.5 can handle both continuous and categorical attributes, while ID3 can only handle categorical attributes. C4.5 uses a method called ``thresholding'' to determine the best split point for continuous attributes, allowing it to create decision trees that can work with a wider variety of data types.
    \end{ejercicio}

    \begin{ejercicio}[Gini Impurity Calculation (3 pts.)]
        Consider the dataset with 20 points given in Table~\ref{tab:fruit_dataset2} and illustrated in Figure~\ref{fig:fruit_scatter_plot2}. Features $F_1$ and $F_2$ correspond to $Feature1$ and $Feature2$. Calculate the individual Gini impurities and the total Gini impurity for the split at $F_2 = 4$ based on the target class \textit{Type} (Pear/Apple).

        \begin{figure}[h!]
        \centering
        \begin{tikzpicture}[scale=0.9, font=\sffamily\scriptsize]
            % Grid
            \draw[step=0.5, gray!20, thin] (0,0) grid (9,6.5);
            \draw[step=1.0, gray!40, stroke] (0,0) grid (9,6.5);
            
            % Axes
            \draw[-Stealth, thick] (0,0) -- (9.2,0) node[below] {Feature1};
            \draw[-Stealth, thick] (0,0) -- (0,6.7) node[above] {Feature2};
            
            % Ticks
            \foreach \x in {1,...,9} \draw (\x,0.05) -- (\x,-0.05) node[below] {\x};
            \foreach \y in {1,...,6} \draw (0.05,\y) -- (-0.05,\y) node[left] {\y};
            
            % Data Points (Pears: Red Squares, Apples: Blue Circles)
            % Pears
            \node[circle, fill=red!80!black, inner sep=1.5pt, label={above:\tiny 1}] at (2.00, 5.75) {};
            \node[circle, fill=red!80!black, inner sep=1.5pt, label={above:\tiny 2}] at (1.25, 3.75) {};
            \node[circle, fill=red!80!black, inner sep=1.5pt, label={above:\tiny 3}] at (3.50, 3.75) {};
            \node[circle, fill=red!80!black, inner sep=1.5pt, label={above:\tiny 4}] at (4.00, 5.75) {};
            \node[circle, fill=red!80!black, inner sep=1.5pt, label={above:\tiny 5}] at (1.00, 2.00) {};
            \node[circle, fill=red!80!black, inner sep=1.5pt, label={above:\tiny 6}] at (2.75, 2.25) {};
            \node[circle, fill=red!80!black, inner sep=1.5pt, label={above:\tiny 7}] at (1.25, 1.25) {};
            \node[circle, fill=red!80!black, inner sep=1.5pt, label={above:\tiny 11}] at (5.75, 5.50) {};
            \node[circle, fill=red!80!black, inner sep=1.5pt, label={above:\tiny 13}] at (7.75, 5.50) {};
            \node[circle, fill=red!80!black, inner sep=1.5pt, label={above:\tiny 15}] at (8.00, 3.50) {};

            % Apples
            \node[circle, fill=blue!80!black, inner sep=1.5pt, label={above:\tiny 8}] at (3.50, 1.75) {};
            \node[circle, fill=blue!80!black, inner sep=1.5pt, label={above:\tiny 9}] at (4.75, 0.50) {};
            \node[circle, fill=blue!80!black, inner sep=1.5pt, label={above:\tiny 10}] at (5.25, 4.25) {};
            \node[circle, fill=blue!80!black, inner sep=1.5pt, label={above:\tiny 12}] at (5.50, 2.50) {};
            \node[circle, fill=blue!80!black, inner sep=1.5pt, label={above:\tiny 14}] at (8.25, 4.50) {};
            \node[circle, fill=blue!80!black, inner sep=1.5pt, label={above:\tiny 16}] at (7.75, 3.75) {};
            \node[circle, fill=blue!80!black, inner sep=1.5pt, label={above:\tiny 17}] at (6.50, 3.50) {};
            \node[circle, fill=blue!80!black, inner sep=1.5pt, label={above:\tiny 18}] at (6.50, 2.25) {};
            \node[circle, fill=blue!80!black, inner sep=1.5pt, label={above:\tiny 19}] at (6.75, 3.25) {};
            \node[circle, fill=blue!80!black, inner sep=1.5pt, label={above:\tiny 20}] at (7.00, 2.00) {};
            
            % Legend
            \begin{scope}[shift={(6.5, 0.8)}]
                \draw[fill=white, draw=gray] (-0.2,-0.2) rectangle (2.2, 0.8);
                \node[circle, fill=red!80!black, inner sep=1.5pt, label={right:\tiny Pear}] at (0,0.5) {};
                \node[circle, fill=blue!80!black, inner sep=1.5pt, label={right:\tiny Apple}] at (0,0.1) {};
            \end{scope}
        \end{tikzpicture}
        \caption{Scatter Plot of Feature1 vs Feature2}
        \label{fig:fruit_scatter_plot2}
        \end{figure}

        \begin{table}[h!]
        \centering
        \small
        \begin{tabular}{|c|c|c|c|}
        \hline
        \textbf{Index} & $F_1$ & $F_2$ & \textbf{Type} \\ \hline
        1 & 2.00 & 5.75 & Pear \\ \hline
        2 & 1.25 & 3.75 & Pear \\ \hline
        3 & 3.50 & 3.75 & Pear \\ \hline
        4 & 4.00 & 5.75 & Pear \\ \hline
        5 & 1.00 & 2.00 & Pear \\ \hline
        6 & 2.75 & 2.25 & Pear \\ \hline
        7 & 1.25 & 1.25 & Pear \\ \hline
        8 & 3.50 & 1.75 & Apple \\ \hline
        9 & 4.75 & 0.50 & Apple \\ \hline
        10 & 5.25 & 4.25 & Apple \\ \hline
        11 & 5.75 & 5.50 & Pear \\ \hline
        12 & 5.50 & 2.50 & Apple \\ \hline
        13 & 7.75 & 5.50 & Pear \\ \hline
        14 & 8.25 & 4.50 & Apple \\ \hline
        15 & 8.00 & 3.50 & Pear \\ \hline
        16 & 7.75 & 3.75 & Apple \\ \hline
        17 & 6.50 & 3.50 & Apple \\ \hline
        18 & 6.50 & 2.25 & Apple \\ \hline
        19 & 6.75 & 3.25 & Apple \\ \hline
        20 & 7.00 & 2.00 & Apple \\ \hline
        \end{tabular}
        \caption{Fruit Dataset}
        \label{tab:fruit_dataset2}
        \end{table}

        When splitting the dataset at $F_2 = 4$, we have two subsets:
        \begin{itemize}
            \item \textbf{Subset 1 ($F_2 < 4$):} The dataset here contains 6 pears (indices 2, 3, 5, 6, 7, 15) and 8 apples (indices 8, 9, 12, 16, 17, 18, 19, 20). The Gini impurity for this subset is calculated as:
            \begin{equation*}
                Gini_1 = 1 - \left(\frac{6}{14}\right)^2 - \left(\frac{8}{14}\right)^2 = \dfrac{24}{49} \approx 0.4898
            \end{equation*}

            \item \textbf{Subset 2 ($F_2 \geq 4$):} The dataset here contains 4 pears (indices 1, 4, 11, 13) and 2 apples (indices 10, 14). The Gini impurity for this subset is calculated as:
            \begin{equation*}
                Gini_2 = 1 - \left(\frac{4}{6}\right)^2 - \left(\frac{2}{6}\right)^2 = \dfrac{4}{9} \approx 0.4444
            \end{equation*}
        \end{itemize}

        The total Gini impurity for the split at $F_2 = 4$ is calculated as a weighted average of the individual Gini impurities:
        \begin{equation*}
            Gini_{total} = \frac{14}{20} \cdot Gini_1 + \frac{6}{20} \cdot Gini_2 = \frac{14}{20} \cdot \dfrac{24}{49} + \frac{6}{20} \cdot \dfrac{4}{9} = \dfrac{10}{21} \approx 0.4762
        \end{equation*}
    \end{ejercicio}

    \begin{ejercicio}[Decision Tree Construction (5 pts.)]
        Table~\ref{tab:feature_splits2} provides several candidate splits and their corresponding Gini Impurities. Create a Decision Tree that correctly classifies all points based on this table. Draw it explicitly and specify its depth.

        \begin{table}[h!]
        \centering
        \begin{tabular}{|c|c|c|}
        \hline
        \textbf{Feature} & \textbf{Split} & \textbf{Gini Impurity} \\ \hline
        $F_1$ & $<1$ & 0.487 \\ \hline
        $F_1$ & $\le 2$ & 0.469 \\ \hline
        $F_1$ & $<3$ & 0.313 \\ \hline
        $F_1$ & $<4$ & 0.444 \\ \hline
        $F_1$ & $<5$ & 0.374 \\ \hline
        $F_1$ & $<6$ & 0.487 \\ \hline
        $F_1$ & $\le 7$ & 0.500 \\ \hline
        $F_1$ & $<8$ & 0.421 \\ \hline
        $F_2$ & $<1$ & 0.498 \\ \hline
        $F_2$ & $<2$ & 0.493 \\ \hline
        $F_2$ & $<3$ & 0.479 \\ \hline
        $F_2$ & $<4$ & 0.493 \\ \hline
        $F_2$ & $<5$ & 0.375 \\ \hline
        \end{tabular}
        \caption{Feature Splits with Gini Impurity}
        \label{tab:feature_splits2}
        \end{table}
    \end{ejercicio}

    \begin{ejercicio}[Bootstrap Samples (1 pt.)]
        Why are bootstrap samples used when building a Random Forest?

        Bootstrap samples are used in Random Forest to introduce diversity among the individual decision trees in the ensemble. By randomly sampling the training data with replacement, each tree is trained on a different subset of the data, which helps to reduce overfitting and improve generalization. This randomness allows the Random Forest to capture a wider variety of patterns in the data, leading to more robust and accurate predictions when aggregating the outputs of all trees.
    \end{ejercicio}

    \begin{ejercicio}[Random Forest vs. Decision Tree Classification (5 pts.)]
        Classify the point $(F_1 = 4, F_2 = 5.25)$ using the Random Forest in Figure ~\ref{fig:random_forest_decision_tree2}. State the prediction of each individual tree as well as the aggregated output. Then, classify the same point using the large Decision Tree on the right (draw the decision path).

        \begin{figure}[h!]
        \centering
        \resizebox{\textwidth}{!}{%
        \begin{tikzpicture}[
            node/.style={circle, draw=black, fill=black, text=white, minimum size=0.95cm, inner sep=0pt, font=\scriptsize\bfseries, align=center},
            leaf/.style={circle, draw=black, fill=black, text=white, minimum size=0.85cm, inner sep=0pt, font=\scriptsize\bfseries},
            edge from parent/.style={draw, -Stealth, ultra thick},
            tree label/.style={font=\small\bfseries, yshift=0.35cm}
        ]

        % Small Tree 1
        \begin{scope}[shift={(0,0)}]
            \node[node] at (0,0) {$F_2 < 4$}
                [sibling distance=1.6cm, level distance=1.3cm]
                child { node[leaf] {C2} }
                child { node[leaf] {C3} };
        \end{scope}

        % Small Tree 2
        \begin{scope}[shift={(3.5,0)}]
            \node[node] at (0,0) {$F_1 \ge 4$}
                [sibling distance=1.6cm, level distance=1.3cm]
                child { node[leaf] {C2} }
                child { node[leaf] {C3} };
        \end{scope}

        % Small Tree 3
        \begin{scope}[shift={(0,-3.5)}]
            \node[node] at (0,0) {$F_1 < 3$}
                [sibling distance=1.6cm, level distance=1.3cm]
                child { node[leaf] {C1} }
                child { node[leaf] {C2} };
        \end{scope}

        % Small Tree 4
        \begin{scope}[shift={(3.5,-3.5)}]
            \node[node] at (0,0) {$F_2 > 6$}
                [sibling distance=1.6cm, level distance=1.3cm]
                child { node[leaf] {C1} }
                child { node[leaf] {C2} };
        \end{scope}

        % Large Decision Tree
        \begin{scope}[shift={(10,0)}]
            \node[node] at (0,0) {$F_1 > 3$}
                [sibling distance=3.6cm, level distance=1.3cm]
                child { node[node] {$F_2 > 4$}
                    [sibling distance=1.8cm, level distance=1.3cm]
                    child { node[node] {$F_2 < 5$}
                        [sibling distance=1.2cm, level distance=1.3cm]
                        child { node[leaf] {C1} }
                        child { node[leaf] {C2} }
                    }
                    child { node[leaf] {C4} }
                }
                child { node[node] {$F_2 < 3$}
                    [sibling distance=1.8cm, level distance=1.3cm]
                    child { node[leaf] {C4} }
                    child { node[node] {$F_1 < 4$}
                        [sibling distance=1.2cm, level distance=1.3cm]
                        child { node[leaf] {C3} }
                        child { node[leaf] {C5} }
                    }
                };
        \end{scope}

        \end{tikzpicture}
        }
        \caption{Random Forest \& Decision Tree}
        \label{fig:random_forest_decision_tree2}
        \end{figure}

        In the random forest, there are four trees. The predictions for the point $(F_1 = 4, F_2 = 5.25)$ are as follows:
        \begin{itemize}
            \item Tree 1: C3
            \item Tree 2: C2
            \item Tree 3: C2
            \item Tree 4: C2
        \end{itemize}
        The aggregated output is C2, as it is the most frequent prediction among the four trees.

        In the large decision tree, the decision path for the point $(F_1 = 4, F_2 = 5.25)$ is as follows:
        \begin{enumerate}
            \item Start at the root node: $F_1 > 3$ (True)
            \item Move to the left child: $F_2 > 4$ (True)
            \item Move to the left child: $F_2 < 5$ (False)
            \item Move to the right child: C2
        \end{enumerate}
        Therefore, the classification for the point $(F_1 = 4, F_2 = 5.25)$ using the large decision tree is also C2.
    \end{ejercicio}

    \newpage

    % ------------------------------------
    % PART 4: ENCODING & FORWARD PROPAGATION
    % ------------------------------------

    \section*{Part 4: Encoding \& Forward Propagation (5 pts.)}

    \begin{ejercicio}[Ordinal Encoding (2 pts.)]
        Given a table column with the elements \{\text{``OH''}, \text{``OR''}, \text{``FI''}, \text{``NY''}, \text{``WY''}, \text{``WI''}\}, perform ordinal encoding on these expressions and present the output in a table. What happens if you now want to add \text{``CA''} to the list?

        The ordinal encoding for the given elements can be represented as shown in Table~\ref{tab:ordinal_encoding}.
        \begin{table}[h!]
        \centering
        \begin{tabular}{|c|c|}
        \hline
        \textbf{Element} & \textbf{Ordinal Encoding} \\ \hline
        FI & 1 \\ \hline
        OH & 2 \\ \hline
        OR & 3 \\ \hline
        NY & 4 \\ \hline
        WI & 5 \\ \hline
        WY & 6 \\ \hline
        \end{tabular}
        \caption{Ordinal Encoding of Elements}
        \label{tab:ordinal_encoding}
        \end{table}

        If we want to add ``CA'' to the list, it would have to be assigned the first ordinal value (1), and all other elements would need to be shifted accordingly. This would change the encoding of all existing elements, which can lead to inconsistencies and misinterpretations in the data representation.
    \end{ejercicio}

    \begin{ejercicio}[Forward Propagation Network (2 pts.)]
        Given the small neural network below with input $x = 12$, calculate the output values after applying the activation functions and write them into the corresponding white boxes.

        \begin{figure}[h!]
        \centering
        \resizebox{0.8\textwidth}{!}{%
        \begin{tikzpicture}[
            neuron/.style={circle, draw, fill=black, text=white, minimum size=1.3cm, font=\small\sffamily\bfseries, align=center},
            box/.style={rectangle, draw, minimum size=0.65cm, thick, fill=white},
            edge/.style={-Stealth, thick},
            weight/.style={font=\small, fill=white, inner sep=1.5pt, rounded corners=1pt, sloped}
        ]

        % Nodes Placement
        \node[font=\large\bfseries] (in) at (-1.5,0) {$x = 12$};
        \node[neuron] (Identity) at (1.5,0) {Identity};
        \node[box] (Identity_box) at (2.9,0) {12};

        \node[neuron] (R1) at (6, 2.2) {ReLU};
        \node[box] (R1_box) at (6, 3.3) {48};

        \node[neuron] (R2) at (6, -2.2) {ReLU};
        \node[box] (R2_box) at (6, -3.3) {0};

        \node[neuron] (Sign) at (10.5,0) {Sign};
        \node[box] (Out_box) at (12.2,0) {1};

        % Connections & Weights
        \draw[edge] (in) -- (Identity);
        \draw[edge] (Identity_box) -- node[weight, above] {$w_{1,1} = 4$} (R1);
        \draw[edge] (Identity_box) -- node[weight, below] {$w_{1,2} = -2$} (R2);

        \draw[edge] (R1) -- node[weight, above] {$w_{2,1} = 0.5$} (Sign);
        \draw[edge] (R2) -- node[weight, below] {$w_{2,2} = 0.5$} (Sign);
        \draw[edge] (Sign) -- (Out_box);

        \end{tikzpicture}
        }
        \caption{Small Feedforward Neural Network}
        \label{fig:feedforward_network}
        \end{figure}
    \end{ejercicio}

    \begin{ejercicio}[Error Calculation - MAE \& MSE (1 pt.)]
        Suppose the neural network outputs $2$, and the target value should be $4$. Calculate the Mean Absolute Error (MAE) and Mean Squared Error (MSE).


        The Mean Absolute Error (MAE) is calculated as the absolute difference between the predicted value and the target value:
        \[
        \text{MAE} = |2 - 4| = 2
        \]

        The Mean Squared Error (MSE) is calculated as the square of the difference between the predicted value and the target value:
        \[
        \text{MSE} = (2 - 4)^2 = (-2)^2 = 4
        \]
    \end{ejercicio}

\end{document}